\documentclass[11pt,a4paper]{article}

% Packages
\usepackage[utf8]{inputenc}
\usepackage[T1]{fontenc}
\usepackage[italian]{babel}
\usepackage{geometry}
\usepackage{hyperref}
\usepackage{xcolor}
\usepackage{listings}
\usepackage{tikz}
\usepackage{booktabs}
\usepackage{fancyhdr}
\usepackage{titlesec}
\usepackage{enumitem}
\usepackage{graphicx}
\usepackage{fontspec}

% Page geometry
\geometry{margin=2.5cm}

% Colors
\definecolor{primary}{RGB}{99,102,241}
\definecolor{secondary}{RGB}{139,92,246}
\definecolor{codebg}{RGB}{248,248,252}
\definecolor{darkbg}{RGB}{30,30,40}
\definecolor{textgray}{RGB}{100,100,120}

% TikZ libraries
\usetikzlibrary{shapes,arrows,positioning,calc,fit,backgrounds,shadows}

% Listings configuration
\lstset{
    backgroundcolor=\color{codebg},
    basicstyle=\small\ttfamily,
    breakatwhitespace=false,
    breaklines=true,
    captionpos=b,
    commentstyle=\color{textgray},
    frame=single,
    keepspaces=true,
    keywordstyle=\color{primary}\bfseries,
    language=JavaScript,
    numbers=left,
    numbersep=5pt,
    showspaces=false,
    showstringspaces=false,
    showtabs=false,
    tabsize=2,
    title=\lstname,
    xleftmargin=15pt,
    framexleftmargin=15pt
}

% Headers
\pagestyle{fancy}
\fancyhf{}
\fancyhead[L]{\textcolor{textgray}{Silvi - PDF Viewer + AI Chat}}
\fancyhead[R]{\textcolor{textgray}{\thepage}}
\fancyfoot[C]{\textcolor{textgray}{Documentazione Tecnica}}
\renewcommand{\headrulewidth}{0.5pt}
\renewcommand{\footrulewidth}{0.3pt}

% Section styling
\titleformat{\section}{\Large\bfseries\color{primary}}{\thesection}{1em}{}
\titleformat{\subsection}{\large\bfseries\color{secondary}}{\thesubsection}{1em}{}
\titleformat{\subsubsection}{\normalsize\bfseries\color{textgray}}{\thesubsubsection}{1em}{}

% Document info
\title{
    \vspace{-2cm}
    \begin{tikzpicture}
        \node[rectangle, rounded corners=10pt, fill=darkbg, minimum width=0.9\textwidth, minimum height=3cm] (box) {};
        \node[below=0.3cm of box.north, text=white, font=\Huge\bfseries] {PDF Viewer + AI Chat};
        \node[below=1.5cm of box.north, text=white!70, font=\large] {Visualizzazione PDF e Tutoring AI con RAG};
    \end{tikzpicture}
    \vspace{1cm}
}
\author{Silvi Documentation Team}
\date{Febbraio 2026}

\begin{document}

\maketitle
\thispagestyle{empty}

\vspace{2cm}

\tableofcontents
\newpage

%==============================================================================
\section{Introduzione}
%==============================================================================

\subsection{Panoramica}

Il componente \textbf{PDF Viewer + AI Chat} implementa un sistema completo di visualizzazione PDF integrato con un \textbf{tutoring AI contestuale} basato sulla tecnologia RAG (Retrieval Augmented Generation).

\subsubsection{Funzionalità Principali}

\begin{itemize}[leftmargin=*]
    \item \textbf{PDF Viewer}: Due implementazioni ottimizzate per diversi casi d'uso
    \item \textbf{AI Tutoring}: Chat contestuale basata sul contenuto del PDF (RAG)
    \item \textbf{Vector Search}: Ricerca semantica con embeddings (Pinecone + OpenAI)
    \item \textbf{Source Grounding}: Tracciamento della fonte per ogni risposta
    \item \textbf{Performance}: Caching LRU per il parsing PDF
\end{itemize}

\subsection{Stack Tecnologico}

\begin{table}[h]
\centering
\begin{tabular}{ll}
\toprule
\textbf{Componente} & \textbf{Tecnologia} \\
\midrule
PDF Rendering & \texttt{react-pdf} + \texttt{@react-pdf-viewer} \\
Vector DB & Pinecone \\
Embeddings & OpenAI \texttt{text-embedding-3-small} \\
Text Splitting & LangChain \texttt{RecursiveCharacterTextSplitter} \\
AI Model & OpenAI GPT-4o-mini \\
Caching & In-memory LRU con TTL \\
\bottomrule
\end{tabular}
\caption{Stack tecnologico del componente PDF + AI}
\end{table}

\subsection{Architettura Generale}

\begin{center}
\begin{tikzpicture}[
    node distance=0.8cm,
    box/.style={rectangle, rounded corners, draw=primary, fill=primary!10, minimum width=2.8cm, minimum height=0.8cm, font=\small},
    smallbox/.style={rectangle, rounded corners, draw=secondary, fill=secondary!10, minimum width=2.5cm, minimum height=0.6cm, font=\scriptsize},
    arrow/.style={->, >=stealth, thick, primary},
    label/.style={font=\tiny, text=textgray}
]

% Frontend column
\node[box] (frontend) {Frontend};
\node[smallbox, below=0.4cm of frontend] (viewer) {PDFViewer};
\node[smallbox, below=0.3cm of viewer] (chat) {PDFChat};

% Backend column
\node[box, right=2.5cm of frontend] (backend) {Backend};
\node[smallbox, below=0.4cm of backend] (express) {Express};
\node[smallbox, below=0.3cm of express] (asktutor) {askTutor()};

% Services
\node[box, right=2.5cm of backend] (services) {Services};
\node[smallbox, below=0.4cm of services] (cache) {PDF Cache};
\node[smallbox, below=0.3cm of cache] (vector) {Vector Store};

% External
\node[box, right=2.5cm of services] (external) {External};
\node[smallbox, below=0.4cm of external] (pinecone) {Pinecone};
\node[smallbox, below=0.3cm of pinecone] (openai) {OpenAI};

% Arrows
\draw[arrow, <->] (viewer.east) -- node[above, label] {PDF URL} (express.west);
\draw[arrow, <->] (chat.east) -- node[below, label] {/api/chat} (asktutor.west);
\draw[arrow, ->] (asktutor.east) -- node[above, label] {Parse} (cache.west);
\draw[arrow, ->] (cache.east) -- node[above, label] {Query} (vector.west);
\draw[arrow, ->] (vector.east) -- node[above, label] {Search} (pinecone.west);
\draw[arrow, ->] (asktutor.north) |- ++(0,0.5) -| node[above, label] {RAG} (openai.north);

\end{tikzpicture}
\end{center}

%==============================================================================
\section{Architettura RAG}
%==============================================================================

\subsection{Retrieval Augmented Generation (RAG)}

Il sistema RAG combina \textbf{retrieval} di documenti con \textbf{generazione} testuale:

\begin{enumerate}
    \item \textbf{Ingestion}: Il PDF viene caricato, parsato e diviso in chunk
    \item \textbf{Vectorization}: Ogni chunk viene convertito in embedding
    \item \textbf{Storage}: Gli embeddings vengono salvati in Pinecone
    \item \textbf{Query}: La domanda dell'utente viene convertita in embedding
    \item \textbf{Retrieval}: Si cercano i chunk più simili nel vector store
    \item \textbf{Generation}: Il contesto + domanda viene inviato a GPT-4o-mini
\end{enumerate}

\begin{center}
\begin{tikzpicture}[
    node distance=0.6cm,
    phase/.style={rectangle, rounded corners=5pt, draw=primary, fill=primary!5, minimum width=2.2cm, minimum height=0.6cm, font=\scriptsize, align=center},
    arrow/.style={->, >=stealth, thick, primary}
]

% Ingestion row
\node[phase] (upload) {PDF Upload};
\node[phase, right=0.4cm of upload] (parse) {Parse\\(pdf-parse)};
\node[phase, right=0.4cm of parse] (split) {Split\\(1000 chars)};
\node[phase, right=0.4cm of split] (embed) {Embed\\(OpenAI)};
\node[phase, right=0.4cm of embed] (store) {Store\\(Pinecone)};

% Query row
\node[phase, below=1.2cm of upload] (question) {Question};
\node[phase, right=0.4cm of question] (qembed) {Embed\\(OpenAI)};
\node[phase, right=0.4cm of qembed] (search) {Search\\(Pinecone)};
\node[phase, right=0.4cm of search] (context) {Context\\(topK=5)};
\node[phase, right=0.4cm of context] (generate) {Generate\\(GPT-4o)};

% Arrows
\draw[arrow] (upload) -- (parse);
\draw[arrow] (parse) -- (split);
\draw[arrow] (split) -- (embed);
\draw[arrow] (embed) -- (store);

\draw[arrow] (question) -- (qembed);
\draw[arrow] (qembed) -- (search);
\draw[arrow] (search) -- (context);
\draw[arrow] (context) -- (generate);

% Dotted connection
\draw[dashed, gray, thick] (store) -- (search);

\end{tikzpicture}
\end{center}

%==============================================================================
\section{PDF Viewer Frontend}
%==============================================================================

\subsection{FluidPDFViewer (react-pdf)}

Visualizzatore PDF ``fluido'' con \textbf{lazy loading} delle pagine e \textbf{responsive} width.

\subsubsection{Caratteristiche}

\begin{itemize}[leftmargin=*]
    \item \textbf{Lazy Loading}: Solo le pagine visibili vengono renderizzate
    \item \textbf{Responsive}: Adattamento automatico alla larghezza del contenitore
    \item \textbf{Worker CDN}: Configurazione worker PDF.js da CDN per evitare problemi MIME/CORS
    \item \textbf{Error Classification}: Categorizzazione errori per messaggi user-friendly
\end{itemize}

\begin{lstlisting}[language=JavaScript, caption={Lazy Loading Logic}]
const usePageVisibility = (numPages, containerRef) => {
  const [visiblePages, setVisiblePages] = useState(new Set());

  useEffect(() => {
    // Carica le prime 5 pagine immediatamente
    const initialPages = new Set();
    for (let i = 1; i <= Math.min(5, numPages); i++) {
      initialPages.add(i);
    }
    setVisiblePages(initialPages);

    const handleScroll = () => {
      const scrollTop = containerRef.current.scrollTop;
      const currentPage = Math.floor(scrollTop / 1000) + 1;
      
      setVisiblePages(prev => {
        const newVisiblePages = new Set(prev);
        // Pre-loading: corrente + 2 precedenti + 5 successive
        for (let i = Math.max(1, currentPage - 2); 
             i <= Math.min(numPages, currentPage + 5); i++) {
          newVisiblePages.add(i);
        }
        return newVisiblePages;
      });
    };

    // Throttle con requestAnimationFrame
    containerRef.current.addEventListener('scroll', handleScroll);
  }, [numPages]);

  return visiblePages;
};
\end{lstlisting}

\subsection{PDFReader (@react-pdf-viewer)}

Visualizzatore PDF avanzato con funzionalità professionali.

\subsubsection{Caratteristiche}

\begin{itemize}[leftmargin=*]
    \item \textbf{Plugin Architecture}: defaultLayout, pageNavigation, search
    \item \textbf{Theme Adaptation}: Sincronizzazione dark/light mode via MutationObserver
    \item \textbf{Imperative API}: \texttt{jumpToPage}, \texttt{highlightText}, \texttt{jumpToPageAndHighlight}
    \item \textbf{Worker Fallback}: Switch da locale a CDN worker su errori MIME/CORS
\end{itemize}

\begin{lstlisting}[language=JavaScript, caption={Imperative API Ref}]
export interface PDFReaderRef {
  jumpToPage: (pageIndex: number) => void;
  highlightText: (text: string, options?: SearchOptions) => void;
  jumpToPageAndHighlight: (pageNumber: number, text: string) => void;
}

// Expose methods
useImperativeHandle(ref, () => ({
  jumpToPage: (pageIndex) => {
    pageNavigationPluginInstance.jumpToPage(pageIndex);
  },
  highlightText: (text, options) => {
    searchPluginInstance.highlight(text, options);
  },
  jumpToPageAndHighlight: (pageNumber, text) => {
    pageNavigationPluginInstance.jumpToPage(pageNumber - 1);
    setTimeout(() => {
      searchPluginInstance.highlight(text, { caseSensitive: false });
    }, 300);
  },
}));
\end{lstlisting}

%==============================================================================
\section{AI Chat Component}
%==============================================================================

\subsection{PDFChat}

Componente chat che utilizza RAG per rispondere basandosi sul PDF caricato.

\subsubsection{Key Features}

\begin{itemize}[leftmargin=*]
    \item \textbf{Context Window}: Mantiene ultimi 12 messaggi come history
    \item \textbf{Grounding Constraint}: System prompt restringe risposte al contenuto PDF
    \item \textbf{Rate Limiting}: Gestita dal backend (\texttt{aiLimiter})
    \item \textbf{Error Handling}: Visualizzazione messaggi errore user-friendly
\end{itemize}

\begin{lstlisting}[language=JavaScript, caption={Send Message con RAG}]
const sendMessage = useCallback(async (raw: string) => {
  const content = raw.trim();
  
  // Aggiungi messaggio utente
  setMessages(prev => [...prev, { role: 'user', content }]);
  setIsSending(true);

  try {
    // Prepara history (ultimi 12 messaggi)
    const history = messagesRef.current
      .filter(m => m.role === 'user' || m.role === 'assistant')
      .slice(-12);

    // Chiama API RAG
    const reply = await studyService.askTutor(deckId, content, history);
    
    // Aggiungi risposta AI
    setMessages(prev => [...prev, { role: 'assistant', content: reply }]);
  } catch (err) {
    setError(err?.message || 'Errore nella chat');
  } finally {
    setIsSending(false);
  }
}, [deckId]);
\end{lstlisting}

%==============================================================================
\section{Vector Store Service}
%==============================================================================

\subsection{Pinecone + LangChain Integration}

Servizio per la gestione degli embeddings e la ricerca vettoriale.

\begin{lstlisting}[language=JavaScript, caption={Vector Store Service}]
const DEFAULT_EMBED_MODEL = 'text-embedding-3-small';
const DEFAULT_TOP_K = 5;

// Configurazione Text Splitter
const getTextSplitter = async () => {
  const { RecursiveCharacterTextSplitter } = 
    await import('langchain/text_splitter');
  return new RecursiveCharacterTextSplitter({
    chunkSize: 1000,
    chunkOverlap: 200,
  });
};

// Ingestion: PDF -> Chunks -> Embeddings -> Pinecone
const ingestDeck = async (deckId, text) => {
  const splitter = await getTextSplitter();
  const chunks = await splitter.splitText(text);
  
  const [index, embeddings] = await Promise.all([
    getPineconeIndex(), 
    getEmbedder()
  ]);

  const vectors = [];
  for (let i = 0; i < chunks.length; i++) {
    const values = await embeddings.embedQuery(chunks[i]);
    vectors.push({
      id: `${deckId}-${Date.now()}-${i}`,
      values,
      metadata: { deckId, text: chunks[i], chunkIndex: i },
    });
  }

  await index.upsert(vectors);
  return { deckId, vectors: vectors.length };
};

// Query: Question -> Embedding -> Similarity Search
const queryDeck = async (deckId, question, topK = DEFAULT_TOP_K) => {
  const [index, embeddings] = await Promise.all([
    getPineconeIndex(), 
    getEmbedder()
  ]);
  
  const queryVector = await embeddings.embedQuery(question);
  
  const result = await index.query({
    vector: queryVector,
    topK,
    includeMetadata: true,
    filter: { deckId },  // Filtro per deck specifico
  });

  return result.matches
    .map(m => m.metadata?.text)
    .filter(Boolean);
};
\end{lstlisting}

%==============================================================================
\section{PDF Cache Service}
%==============================================================================

\subsection{LRU Cache con TTL}

Servizio di caching in-memory per il parsing PDF.

\subsubsection{Caratteristiche}

\begin{itemize}[leftmargin=*]
    \item \textbf{Chiave Cache}: Hash MD5 del file
    \item \textbf{TTL}: 5 minuti
    \item \textbf{Max Size}: 50 PDF
    \item \textbf{Eviction}: LRU (Least Recently Used)
    \item \textbf{Stats}: Hits, misses, evictions per monitoring
\end{itemize}

\begin{lstlisting}[language=JavaScript, caption={PDF Cache Service}]
class PDFCacheService {
  constructor() {
    this.cache = new Map();
    this.MAX_CACHE_SIZE = 50;
    this.TTL = 5 * 60 * 1000; // 5 minuti
    
    this.stats = { hits: 0, misses: 0, evictions: 0 };
  }

  async parsePDF(filePath, buffer = null) {
    // Calcola hash per cache key
    const hash = buffer 
      ? crypto.createHash('md5').update(buffer).digest('hex')
      : await this._getFileHash(filePath);
    
    const cacheKey = `pdf_${hash}`;
    
    // Check cache
    const cached = this.cache.get(cacheKey);
    if (cached) {
      this.stats.hits++;
      this._touchEntry(cacheKey);
      return cached.data;
    }
    
    // Cache MISS: parsa PDF
    this.stats.misses++;
    const pdfData = await pdfParse(buffer || await fs.readFile(filePath));
    
    // Evict se necessario
    this._evictIfNeeded();
    
    // Salva in cache
    this.cache.set(cacheKey, {
      data: pdfData,
      timestamp: Date.now(),
    });
    
    return pdfData;
  }

  _evictIfNeeded() {
    if (this.cache.size >= this.MAX_CACHE_SIZE) {
      const entries = [...this.cache.entries()];
      const oldest = entries.sort((a, b) => 
        a[1].timestamp - b[1].timestamp)[0];
      
      this.cache.delete(oldest[0]);
      this.stats.evictions++;
    }
  }
}
\end{lstlisting}

%==============================================================================
\section{Backend Integration}
%==============================================================================

\subsection{askTutor Service}

Servizio backend che implementa il RAG per la chat AI.

\begin{lstlisting}[language=JavaScript, caption={askTutor - RAG Implementation}]
async askTutor(tenantScope, deckId, message, history = []) {
  // 1. Validazioni
  if (!message || message.length > 2000) {
    throw AppError.validation('Messaggio non valido');
  }

  // 2. Recupera deck e verifica testo estratto
  const deck = await Deck.findOne({ _id: deckId, user: userId });
  if (!deck.extractedText || deck.extractedText.length < 100) {
    throw AppError.validation('Testo PDF non disponibile');
  }

  // 3. QUERY VETTORIALE: Trova chunk rilevanti
  const relevantChunks = await vectorStoreService.queryDeck(
    deckId, message.trim(), 5
  );
  
  if (!relevantChunks || relevantChunks.length === 0) {
    return "Non ho trovato informazioni rilevanti nel PDF...";
  }

  // 4. ASSEMBLY CONTEXTO
  const context = relevantChunks.join('\n\n---\n\n');
  const truncatedContext = context.length > 4000 
    ? context.substring(0, 4000) + '...' 
    : context;

  // 5. COSTRUZIONE MESSAGGI
  const systemPrompt = `Sei un tutor AI che risponde SOLO basandosi 
sul contenuto fornito.
REGOLE:
1. Rispondi SOLO usando le informazioni nel contesto
2. Se l'info non e' nel contesto, dillo chiaramente
3. Non fare supposizioni oltre il testo fornito`;

  const messages = [
    { role: 'system', content: systemPrompt },
    ...history.slice(-12),
    { role: 'user', content: `CONTESTO: ${truncatedContext}\n\nDOMANDA: ${message}` }
  ];

  // 6. CHIAMATA OPENAI
  const completion = await openai.chat.completions.create({
    model: 'gpt-4o-mini',
    messages,
    max_tokens: 1000,
    temperature: 0.3,  // Bassa temperatura per factual
  });

  return completion.choices[0].message.content.trim();
}
\end{lstlisting}

%==============================================================================
\section{API Endpoints}
%==============================================================================

\begin{table}[h]
\centering
\begin{tabular}{llp{6cm}}
\toprule
\textbf{Endpoint} & \textbf{Method} & \textbf{Descrizione} \\
\midrule
\texttt{/api/study/:id/generate-pdf} & POST & Upload PDF + generazione flashcard AI \\
\texttt{/api/study/:id/chat} & POST & Chat AI Tutor con RAG \\
\texttt{/api/study/:id/answer-question} & POST & Risposta a domanda d'esame (contesto PDF) \\
\bottomrule
\end{tabular}
\caption{API Endpoints per PDF + AI}
\end{table}

%==============================================================================
\section{Glossario}
%==============================================================================

\begin{table}[h]
\centering
\begin{tabular}{lp{9cm}}
\toprule
\textbf{Termine} & \textbf{Definizione} \\
\midrule
\textbf{RAG} & Retrieval Augmented Generation - tecnica che combina retrieval di documenti con generazione testuale \\
\textbf{Embedding} & Rappresentazione vettoriale di un testo che cattura il significato semantico \\
\textbf{Vector Store} & Database specializzato per la ricerca per similarità tra vettori \\
\textbf{Chunking} & Divisione di un testo lungo in segmenti più piccoli per l'elaborazione \\
\textbf{Similarity Search} & Ricerca dei vettori più simili a una query (coseno similarity) \\
\textbf{Source Grounding} & Ancoraggio della risposta AI al documento fonte \\
\textbf{LRU Cache} & Least Recently Used - strategia di caching che elimina prima gli elementi meno usati \\
\textbf{TTL} & Time To Live - durata di vita di un elemento in cache \\
\bottomrule
\end{tabular}
\caption{Glossario dei termini tecnici}
\end{table}

%==============================================================================
\section{Configurazione Environment}
%==============================================================================

\begin{lstlisting}[language=bash, caption={Variabili d'ambiente richieste}]
# Vector Store (Pinecone)
PINECONE_API_KEY=pc_xxx...
PINECONE_INDEX=silvi-vectors

# Embeddings (OpenAI)
OPENAI_API_KEY=sk-...
OPENAI_EMBED_MODEL=text-embedding-3-small

# AI Model
OPENAI_MODEL=gpt-4o-mini
\end{lstlisting}

\vfill

\begin{center}
\textcolor{textgray}{\rule{0.5\textwidth}{0.5pt}}

\textit{Documento generato automaticamente da Kimi Code CLI}

\textcolor{textgray}{Febbraio 2026}
\end{center}

\end{document}
